\chapter{Spectral Rigidity}

\section*{Standing Assumptions}

Throughout this chapter:

\begin{itemize}
\item $\mathcal{H}$ is a separable Hilbert space.
\item $A : \mathcal{D}(A) \subset \mathcal{H} \to \mathcal{H}$ is densely defined, self-adjoint.
\item $A$ is nonnegative: $\langle Ax,x\rangle \ge 0$ for all $x\in\mathcal{D}(A)$.
\item $\ker A$ is finite-dimensional.
\end{itemize}

\section{Rayleigh Quotient}

\begin{definition}[Rayleigh Quotient]
For $x \in \mathcal{D}(A)\setminus\{0\}$,
\[
R_A(x) := \frac{\langle Ax,x\rangle}{\|x\|^2}.
\]
\end{definition}

\section{Spectral Gap}

\begin{definition}[Spectral Gap]
The spectral gap of $A$ is
\[
\lambda_1(A)
:=
\inf_{\substack{x \in \mathcal{D}(A)\\ x \perp \ker A \\ \|x\|=1}}
\langle Ax,x\rangle.
\]
\end{definition}

\begin{lemma}
If $\lambda_1(A)>0$, then
\[
\langle Ax,x\rangle \ge \lambda_1(A)\|x\|^2
\quad
\text{for all } x \perp \ker A.
\]
\end{lemma}

\begin{proof}
Immediate from the definition of $\lambda_1(A)$.
\end{proof}

\section{Min--Max Principle}

\begin{theorem}[Variational Characterization]
Under the standing assumptions,
\[
\lambda_1(A)
=
\inf_{\substack{x \perp \ker A \\ \|x\|=1}}
\langle Ax,x\rangle.
\]
\end{theorem}

\begin{proof}
Follows from the standard min--max principle for self-adjoint operators.
\end{proof}

\section{Operator Domination}

\begin{definition}[Domination]
Let $A,B$ be self-adjoint on $\mathcal{H}$.  
We say $A \succeq B$ if
\[
\langle Ax,x\rangle \ge \langle Bx,x\rangle
\quad
\forall x \in \mathcal{D}(A)\cap\mathcal{D}(B).
\]
\end{definition}

\begin{theorem}[Gap Preservation Under Domination]
If $A \succeq B$ and $\lambda_1(B)\ge \mu>0$, then
\[
\lambda_1(A)\ge \mu.
\]
\end{theorem}

\begin{proof}
For $x \perp \ker B$,
\[
\langle Ax,x\rangle \ge \langle Bx,x\rangle
\ge \mu \|x\|^2.
\]
Taking the infimum yields the claim.
\end{proof}

\section{Local-to-Global Coercivity}

\begin{theorem}[Finite Orthogonal Decomposition]
Assume
\[
\mathcal{H} = \bigoplus_{i=1}^m \mathcal{H}_i
\]
orthogonally and
\[
A = \sum_{i=1}^m A_i
\]
with each $A_i$ nonnegative and coercive:
\[
\langle A_i x_i,x_i\rangle \ge c_i \|x_i\|^2.
\]
Then
\[
\langle Ax,x\rangle
\ge
\left(\min_i c_i\right)\|x\|^2.
\]
\end{theorem}

\begin{proof}
Write $x=\sum_i x_i$. Then
\[
\langle Ax,x\rangle
=
\sum_i \langle A_i x_i,x_i\rangle
\ge
\sum_i c_i \|x_i\|^2
\ge
\left(\min_i c_i\right)\sum_i \|x_i\|^2.
\]
\end{proof}

\section{Strong Resolvent Stability}

\begin{theorem}
Let $A_n \to A$ in strong resolvent sense and assume
\[
\inf_n \lambda_1(A_n) \ge \lambda_* > 0.
\]
Then
\[
\lambda_1(A) \ge \lambda_*.
\]
\end{theorem}

\begin{proof}
Lower semicontinuity of the spectrum under strong resolvent convergence preserves the lower bound.
\end{proof}
